\vspace{0.5em}
\raggedright
\footnotesize
\textit{Notes:} Difference-in-differences estimates on the estimation sample of Table~\ref{tab:did_results}, log points unless stated; means across markets. Panel~B is the comparison the table is built around. The core model predicts a non-disclosure premium of $\hat\beta_1\hat e^*$, which falls short of the observed premium by \attLamApt{} log points in apartments and by \attLamHouse{} in houses. With limited attention the predicted premium is $\hat\beta_1[\hat\psi\hat e^*+(1-\hat\psi)\eta]$, which matches the observed premium in both segments. In apartments the match uses one parameter, $\hat\psi$, estimated from the ratio of the same two aggregates, so it shows that a single attention parameter is enough and is not an out-of-sample test. The evidence for limited attention is in panel~C: attention to silence is a small fraction of attention to ratings in apartments (\attPsiApt{}) and close to full in houses (\attPsiHouse{}), and full attention is rejected for apartments and not for houses. Aggregates weight markets by $\hat\beta_{1,m}$, the weights implied by the ratio-of-sums estimator. $\eta_m$ and $\bar C^{\rm sil}$ value registry letters at DPE band midpoints; values are provisional pending certified consumption. Confidence intervals from a bootstrap over markets (\attNBoot{} draws); percentile intervals are not truncated to the admissible range $[0,1]$ of $\psi_\delta$. The identified interval allows any default $C^{\rm def}\in[0,C^*_\delta]$. The attention ceiling is $\hat\Delta=\eta_m\hat\beta_{1,m}$ (Section~\ref{ss:att_estimates}).
\end{table}
